% !TeX root = ../main.tex

\chapter{绪论}

\section{研究背景及意义}

交通信息工程正在由交通状态监测逐步走向面向事件、对象和场景的综合认知。道路传感器、交通管理系统、地图平台、气象观测和视觉设备持续产生交通数据；这些数据分别描述事故时间、空间位置、道路环境、气象条件、交通设施及事件影响，却以数据表、短文本、多媒体观测和时序记录等形式分散存在。面向交通管理和公众信息服务时，用户通常不会直接操作字段或数据库，而是提出“某起事故发生在何处”“当时记录了什么天气”“哪些事故具有相似道路和环境条件”等自然语言问题。因此，如何将异构交通数据组织为与当前任务相关、能够追溯证据的结构，并准确回答问题，是交通信息智能服务的基础问题。

传统结构化查询可以准确回答字段明确的简单问题，但依赖使用者掌握数据模式与查询语言，对跨实体、跨记录和多条件问题的表达能力有限。通用语言模型具有良好的语言理解与生成能力，却可能在缺少可靠外部证据时生成与数据库不一致的内容。知识图谱能够以实体、关系和属性组织交通事件，将相互独立的数据记录连接为可查询的语义网络；问题驱动的图学习则可从大规模候选知识中选择与当前问题有关的局部结构，降低冗余信息传播。将知识图谱、图神经网络和状态空间模型结合，可形成“数据到知识、知识到证据、证据到答案”的处理链。

US-Accidents 是由 Moosavi 等构建的美国交通事故公开数据集，包含事故位置、时间、自然语言描述、天气、昼夜条件和道路设施等信息。已有研究选取罗德岛州数据构建了以事故为中心的多关系异构知识图谱，为交通事故结构化表示提供了基础。该数据既包含适合图谱映射的类别与布尔字段，也包含可形成文本证据的事故描述，适合检验任务驱动知识图构造与答案预测方法。

本论文以多层次图增强为共同方法主线，围绕交通信息认知与交通场景生成组织两条研究分支。交通信息认知分支由知识问答扩展到视觉与知识结合的场景问答，研究任务驱动图构造、异构语义增强以及问题条件下的区域--属性推理。交通场景生成分支面向多视角时序观测，研究对象关系与时空表征增强。两条分支通过关系构造、表征增强与冗余抑制建立方法联系，分别输出语义答案与四维场景表示。

上述组织方式使知识问答与视觉问答共享交通概念和推理任务，同时保留四维生成在几何表示、观测组织和评价目标上的独立性。已有通用任务上的研究为图增强机制提供方法基础，交通领域中的适用性则需要通过相应的数据构建和迁移实验检验。

\section{国内外研究现状}

\subsection{交通知识图构造与知识问答研究现状}

交通事故研究常使用事故登记、交通事件流、气象、路网、兴趣点和人口统计等数据。早期工作多基于单一事故表开展统计建模，随后逐步引入多源时空数据以分析环境因素、道路条件和时间规律。US-Accidents 通过多类实时交通事件接口收集美国事故记录，并利用反向地理编码、气象、道路设施和昼夜信息进行增强，为大范围交通事故研究提供了统一字段的数据基础。

表格数据能够保留精确数值，却难以直接表达事故之间因共享城市、街道、天气或设施而形成的联系。交通知识图谱通过本体、实体、关系和属性表达多源信息，可支持关联检索和图学习。对于混合型事故数据，连续量全部节点化会产生大量低频节点，类别字段全部保留为属性又会损失共享语义关系。较合理的方式是将事故作为核心实体，将城市、道路、天气、时段和设施等可共享概念构造成节点，同时把经纬度、温度、湿度和风速等连续观测作为事故属性保存。现有交通事故图研究多数服务于严重度分类、风险预测或致因分析，问答任务受到的关注相对有限，仍需与自然语言问题和答案监督结合，形成面向信息服务的问答数据与评价体系。

\subsection{知识问答异构语义表征增强研究现状}

知识问答需要从结构化知识库、文本语料或二者融合的知识源中定位问题相关信息。对于单一事实问题，实体链接与一跳关系查询通常可以获得答案；对于桥接、比较和组合问题，系统需要经过多个中间实体组织证据。固定跳数方法容易在简单问题中引入无关节点，也可能在复杂问题中提前停止。自适应多跳检索通过问题语义、当前检索结果和局部图状态决定下一步扩展，可在检索范围与知识完整性之间取得平衡。

基于图学习的知识问答方法通过上下文相关的多跳检索构建问题相关图，再结合图卷积与选择性过滤预测答案。面向交通事故数据时，此类方法还需结合表格字段语义、时空约束与连续数值特点，组织可供关联推理的知识。

图神经网络通过邻域信息聚合学习实体表示，适合处理交通事故知识图中的共享地点、天气、时段和设施关系。普通图卷积容易把高频公共节点连接的大量信息同时传播给事故实体。问题条件过滤可根据当前问题抑制无关节点和边，使局部图表示更聚焦于回答所需证据。状态空间模型能够以近线性复杂度建模长序列，并通过输入相关参数实现选择性信息保留；将图节点按检索层次组织为序列，再使用问题条件状态空间模块进行过滤，可补充局部图卷积对长程依赖建模的不足。

\subsection{交通场景视觉问答与区域—属性图学习研究现状}

交通场景视觉问答需要结合问题理解图像中的参与者、道路结构、交通控制要素及其关系。目标类别与位置提供基础观测，关系理解与知识推理进一步组织回答所需的信息。已有 PerGraph 工作研究图像质量与美学评价中的区域--属性动态图建模，其动态图交互机制可作为视觉问答扩展的研究基础；问答任务需要进一步引入问题条件和答案监督，并独立验证其有效性。

\subsection{四维交通场景生成与时空关系建模研究现状}

四维交通场景生成需要在多视角、多时刻观测间建立对象级关系，并处理动态目标、遮挡和稀疏视角等问题。跨视角和跨时刻的关联质量直接影响动态场景表示。该分支使用独立的多视角时序输入和几何监督，与问答分支共享关系建模思想，不依赖问答模型的输出。

\section{存在的主要问题}

\subsection{任务相关知识的动态组织问题}

静态事故知识图按照预定义模式连接全部记录，能够表达数据关系，却未必适合具体问题。相同城市、常见天气和高频时段可能连接大量事故，直接在全图传播会引入冗余。需要从问题实体出发动态选择三元组，并根据新增知识与相关性决定检索深度。

\subsection{异构语义的对齐与噪声抑制问题}

事故字段、自然语言描述和派生类别具有不同的数据分布。字段可能缺失，文本可能含有缩写与模板噪声，派生类别也可能损失细节。需要保留来源与缺失状态，区分观测事实和统计插补，并在异构表示增强中抑制冲突信息。

\subsection{问题条件下的视觉区域、属性与关系推理问题}

同一交通图像中的不同问题对应不同的参与者、属性和关系。仅识别目标类别不足以完成关系与知识推理，需要围绕问题动态组织区域--属性关联，并抑制与回答无关的信息。

\subsection{四维交通场景稠密时空证据的组织问题}

多视角时序交通数据中的对象在不同观测下具有位置变化、可见性差异和遮挡关系。需要将局部视觉证据组织为稳定的对象节点与类型化时空关系，抑制错误关联对稠密重建的影响。

\section{论文主要研究内容}

\subsection{面向交通事故知识问答的自适应知识图增强}

面向交通事故知识问答，构建保留数据来源和字段语义的交通事故知识图，研究上下文相关的自适应多跳检索、关系感知图卷积和问题条件状态空间过滤，预测具有事实依据的答案。

\subsection{面向知识问答的异构语义表征增强}

面向交通知识问答中的文本--图知识异构表征，研究分布对齐、语义去噪和非线性融合，降低文本语义与结构化知识之间的表征差异。

\subsection{面向交通场景视觉问答的区域—属性图增强}

以 PerGraph 的区域--属性动态图为方法基础，研究问题条件下的区域与属性激活、跨图交互和一致性约束，形成面向交通场景关系理解与知识推理的视觉问答扩展。原始感知评价实验与交通问答实验分别承担通用机制验证与领域扩展验证。

\subsection{面向四维交通场景的对象关系增强}

面向四维动态交通场景生成，构建潜在对象节点和类型化时空关系，将对象级增强证据回写到稠密场景表示，提升动态区域与遮挡条件下的重建质量。

\section{多层次任务体系与数据关联}

\subsection{交通知识图谱问答数据与任务协议}

第二章将交通数据的实体关系建模、知识图构造与答案预测组织为完整研究内容，并在实验部分介绍问题、答案和支持知识路径的构建。第三章面向异构知识增强问答，根据任务需要组织候选答案与干扰项。两章共享事实来源与交通概念体系，并分别保留适合各自方法的任务协议。

\subsection{交通场景视觉问答数据与图文关联}

第四章围绕问题条件下的区域--属性图推理开展交通场景视觉问答研究，在该章实验部分介绍图像或视频、问题、答案、视觉区域及语义关系的组织方式。知识问答数据与视觉问答数据通过道路、参与者、环境和交通控制等概念建立语义关联；跨数据源的样本对应须以真实采集信息为依据。

\subsection{四维交通场景生成的独立数据体系}

第五章单独介绍多视角时序观测、相机参数及场景重建监督。该数据体系服务于动态场景的几何与外观恢复，与问答数据共享交通应用背景和图增强方法思想，不要求逐样本配对，也不以视觉问答结果作为生成输入。

\section{论文总体技术路线}

交通信息认知分支围绕自然语言交互展开：由任务驱动知识图构造获得问题相关事实，通过异构语义增强改善知识利用，再通过问题条件区域--属性图推理拓展到交通场景视觉问答。各方法章节结合自身研究问题介绍数据构建、任务协议与实验设置。

交通场景生成分支围绕多视角时序建模展开：从视觉观测中组织对象与时空关系，通过关系传播增强场景表征，恢复四维动态交通场景。两条分支共同研究任务相关关系的构造、有效信息的传播及冗余关联的抑制，分别采用问答指标、几何与渲染指标验证方法。

\section{论文章节安排}

第一章介绍研究背景、研究现状、主要问题、研究内容与总体技术路线。第二章研究基于自适应知识图增强的交通事故知识问答，将交通实体关系建模、知识组织和问答样本构建纳入本章。第三章基于 MEAD，研究交通知识问答中的异构语义表征增强与去噪。第四章以 PerGraph 为方法基础，研究问题条件下的区域--属性图增强与交通场景视觉问答，并在本章介绍相应数据及实验协议。第五章独立研究对象关系增强的四维交通场景生成，介绍多视角时序数据及几何与外观评价。第六章总结全文并讨论未来研究方向。
